\section{Results}
\subsection{Datasets and Evaluation Protocol}
\textbf{BCI Competition IV Dataset 4.} We use the released 400-second training recordings and official 200-second test recordings for S1--S3, with 62, 48, and 64 electrodes, respectively~\cite{competition2008}. Final training budgets are 112, 69, and 104 epochs. The supplied full-training checkpoints were verified against their source weights and re-evaluated from raw ECoG using the consolidated implementation. Each prediction contains $200{,}000\times5$ values. The primary score follows the competition's arithmetic mean of twelve correlations: four fingers for each of three subjects. All original test timestamps are scored.

\textbf{Miller fingerflex extension.} We evaluate bp, cc, ht, jc, jp, mv, wc, wm, and zt~\cite{miller2019}, with 38--64 electrodes and approximately 179--610~s per recording. For $N$ original samples, the test boundary is $b=40\lfloor(2N/3)/40\rfloor$. Development uses $[0,b-2000)$ and testing $[b,N)$, with a two-second guard. The development interval's last 20\% is inner validation, separated by another two-second guard. Intervals are processed separately; electrode selection and scaling are refitted in each phase.

Inner training uses at most 100 epochs and patience 25, selecting by validation-interior four-finger correlation. The final budget is that epoch multiplied by the development/inner-training duration ratio, rounded and clipped to 8--160. Final models are trained afresh with common architecture and optimization settings. Test labels remain local; checkpoints and predictions are frozen before local scoring. This follows the competition's chronological split principle but is not an official nine-subject split. BCI recordings originate from this library, so the cohorts overlap.

We calculate each subject/finger Pearson correlation $r_{sf}$ over its complete test interval. With $\mathcal F_4=\{1,2,3,5\}$ and $\mathcal F_5=\{1,2,3,4,5\}$,
\begin{equation}
\bar r_q=\frac{1}{S}\sum_{s=1}^{S}\frac{1}{q}\sum_{f\in\mathcal F_q}r_{sf},\qquad q\in\{4,5\}.
\end{equation}
Subjects receive equal weight. We predict the ring finger but use \rFour{} as the official BCI metric and \rFive{} as the main Miller metric. BCI MAE and MSE average the same four fingers and subjects in original glove units. Correlation measures regression, not classification accuracy.

\begin{table}[t]
\caption{Full test results. BCI: official 200-second tests. Miller: final chronological third with a two-second guard.}
\label{tab:subjects}
\centering
\footnotesize
\begin{tabular}{llrrr}
\toprule
Benchmark & Subject & Test (s) & \rFour & \rFive\\
\midrule
BCI & S1 & 200.00 & 0.8381 & 0.8457\\
 & S2 & 200.00 & 0.6963 & 0.7044\\
 & S3 & 200.00 & 0.7879 & 0.8006\\
 & \textbf{Mean} & & \textbf{0.7741} & \textbf{0.7835}\\
\midrule
Miller & bp & 203.36 & 0.6999 & 0.6993\\
 & cc & 203.36 & 0.7435 & 0.7558\\
 & ht & 203.36 & 0.5218 & 0.5411\\
 & jc & 176.84 & 0.6771 & 0.6633\\
 & jp & 155.28 & 0.6163 & 0.6310\\
 & mv & 59.68 & 0.8812 & 0.8697\\
 & wc & 203.36 & 0.6167 & 0.5916\\
 & wm & 148.28 & 0.2121 & 0.2589\\
 & zt & 203.36 & 0.7676 & 0.7780\\
 & \textbf{Mean} & & \textbf{0.6373} & \textbf{0.6432}\\
\bottomrule
\end{tabular}
\end{table}


\subsection{Performance on the Two Benchmark Settings}
Table~\ref{tab:subjects} gives the complete subject-level results. On BCI, RecurFlex achieves \rFour{} of 0.7741 and \rFive{} of 0.7835 over all 600,000 test samples across subjects. Four-finger MAE and MSE are 0.4294 and 0.4531. All three subjects obtain four-finger correlations above 0.69; the strongest is S1 at 0.8381. Under the original metric, the mean exceeds the historical winning score by approximately 0.3141, without restricting evaluation to active movements or selected test subsequences.

On the nine-subject extension, \rFive{} is 0.6432 and \rFour{} is 0.6373 over 1,556,880 original test samples. The five-finger subject standard deviation is 0.1756; it describes between-subject variability in this single-seed experiment. Performance ranges from 0.2589 for wm to 0.8697 for mv. The little finger has the lowest mean correlation, 0.5396, indicating residual heterogeneity across subjects and fingers.

\subsection{Controlled Frequency-Feature Ablations}
Table~\ref{tab:ablation} reports full official-test scores from fresh feature-variant fits. The first group uses a common 43-slot model, 6,805,126 parameters, identical initialization, electrodes, sampled windows, and budgets. Unused slots are zeroed after scaling. Signed low-frequency input alone obtains \rFour{} of 0.5574, high-frequency power alone 0.7160, and their three-band combination 0.7584. Fusion improves over power alone by 0.0423, with gains for all three subjects in this controlled experiment.

The second group compares three versus four signed bands with a common 44-slot model and matched initialization. Splitting the 4--13-Hz band into theta and alpha improves \rFour{} from 0.7624 to 0.7741, a mean change of 0.0117. Subject changes are $+0.0448$, $+0.0021$, and $-0.0119$; the gain is concentrated in S1 and is not universal. A separate 40-center, linearly spaced 0.5--300-Hz power representation obtains 0.6839, below the high-frequency log-grid representation. Taken together, the experiments support preserving high-frequency spectral resolution and signed low-frequency information. They do not independently establish the benefit of the BiGRU, for which a matched removal ablation remains necessary.

\begin{table}[t]
\caption{Frequency-feature ablations on the complete BCI tests. Parameter-matched comparisons are grouped; $\mathrm{LF}_3$ uses 0.5--4, 4--13, 13--30~Hz.}
\label{tab:ablation}
\centering
\footnotesize
\begin{tabular}{lrr}
\toprule
Representation & \rFour & \rFive\\
\midrule
\multicolumn{3}{l}{\emph{Matched 43-slot experiment}}\\
High-frequency Morlet power (\HF) & 0.7160 & 0.7273\\
Signed low-frequency bands ($\mathrm{LF}_3$) & 0.5574 & 0.5763\\
$\HF+\mathrm{LF}_3$ & 0.7584 & 0.7714\\
\midrule
\multicolumn{3}{l}{\emph{Matched 44-slot band-splitting experiment}}\\
$\HF+\mathrm{LF}_3$; alpha slot zeroed & 0.7624 & 0.7749\\
$\HF+\mathrm{LF}_4$ (RecurFlex) & \textbf{0.7741} & \textbf{0.7835}\\
\midrule
\multicolumn{3}{l}{\emph{Separate frequency-grid comparison}}\\
Linear-grid 0.5--300-Hz Morlet power & 0.6839 & 0.6962\\
\bottomrule
\end{tabular}
\end{table}


\subsection{Comparison with Published Results and Limitations}
Table~\ref{tab:literature} lists BCI scores reported by the TRACE authors~\cite{trace2026}; these baselines were not retrained by us. Our five-finger mean, 0.7835, exceeds their reported TRACE mean of 0.724 by approximately 0.0595, with numerically higher scores for all three subjects. This comparison uses \rFive{}, not the original competition's \rFour{}.

The repository reports Gaussian-smoothed predictions ($\sigma=6$). Its released TRACE configuration shifts targets by 0.2~s; the evaluator scores five outputs at 100~Hz after delay and stride cropping. RecurFlex has no explicit shift or smoothing and scores all original 1,000-Hz timestamps. The released BCI configuration uses the official test as its validation monitor for checkpoint selection. These code-level findings do not verify historical README runs~\cite{trace2026}.

Other BCI reports include HiLoFuseNet's 0.674~\cite{hilofusenet}, switching regression's approximately 0.82~\cite{switching2016}, and BC4D4's approximately 0.85~\cite{bc4d4}. The latter two require reconciliation of scoring and sample selection with our full-original-test evaluation. Unsuccessful reproduction does not invalidate the published BC4D4 result.

On Miller, our 0.6432 is numerically above TRACE's 0.508 and enhanced CORTEG Base's 0.583. TRACE uses an 85/15 split, whereas we test the final third. HiLoFuseNet and CORTEG use the first 400~s for long recordings and 2:1 splits otherwise; their native 25-Hz targets and past-context inputs also differ from our offline reconstruction~\cite{hilofusenet,corteg}. CORTEG's pooled or transfer fitting further differs from our subject-specific training. These are cross-protocol comparisons.

\begin{table}[t]
\caption{Reported five-finger correlations. BCI baselines are TRACE authors' reimplementations; TRACE is their reported run. RecurFlex was evaluated independently. Source means are retained as reported. Miller rows use differing splits and contexts. This is not a common-protocol ranking.}
\label{tab:literature}
\centering
\footnotesize
\begin{tabular}{lrrrr}
\toprule
Method & S1 & S2 & S3 & Mean\\
\midrule
\multicolumn{5}{l}{\emph{BCI: TRACE repository reports}~\cite{trace2026}}\\
FingerFlex & 0.575 & 0.520 & 0.692 & 0.597\\
DTCNet & 0.696 & 0.598 & 0.747 & 0.680\\
DeepFingerNet & 0.427 & 0.312 & 0.541 & 0.427\\
TRACE & 0.776 & 0.622 & 0.773 & 0.724\\
\textbf{RecurFlex} & \textbf{0.8457} & \textbf{0.7044} & \textbf{0.8006} & \textbf{0.7835}\\
\midrule
\multicolumn{5}{l}{\emph{Miller: reported means under differing protocols}}\\
FingerFlex~\cite{fingerflex} & \multicolumn{3}{c}{Original paper} & $\approx$0.49\\
HiLoFuseNet~\cite{hilofusenet} & \multicolumn{3}{c}{Updated manuscript} & 0.558\\
CORTEG~\cite{corteg} & \multicolumn{3}{c}{Pooled} & 0.554\\
CORTEG~\cite{corteg} & \multicolumn{3}{c}{Enhanced Base} & 0.583\\
FingerFlex~\cite{trace2026} & \multicolumn{3}{c}{TRACE authors; 85/15} & 0.436\\
TRACE~\cite{trace2026} & \multicolumn{3}{c}{85/15 split} & 0.508\\
\textbf{RecurFlex} & \multicolumn{3}{c}{Two-thirds split} & \textbf{0.6432}\\
\bottomrule
\end{tabular}
\end{table}


These results support strong offline performance relative to reported state-of-the-art methods; superiority requires matched data, splits, targets, grids, and context. Earlier method exploration inspected the public BCI test, so final prediction freezing does not imply test-blind development. Single-seed results are descriptive. Additional seeds, matched recurrent/context ablations, and unused test data would strengthen generalization evidence.
